\scanpage{98}
\margin{On aura intérêt ici de prendre pour $S$ le mod\unsure{è|e}le complet régulier can. de $K$ sur $k\ldots$}
un corps $k$, alors les $k(s)$ (pour $s\ne\eta$, pt gén.) apparaissent comme des extensions finies de ce $k$, et tout $s\in S$ définit une classe de conjugaison de sous groupes $\unsure{\Delta_{s'}|I_{s'}}$ dans $\mathcal E_K$ (pour sa structure : \unsure{celle|cette} $=$ \inserted{intérieure} \unsure{relation|relative} sur $\mathcal E_k\ldots$), \inserted{et les $N_{s'}$ qui} \deleted{\unsure{qui}} sont justement les sous groupes \deleted{\unsure{normalisateurs}} de \deleted{inertie} \inserted{décomposition} de tantôt, et on peut regarder les sections de l'extension \inserted{de $N_{s'}$ sur} des $\mathcal E_{k(s)}\underset{\inserted{\text{ouvert}}}{\subset}\mathcal E_K$ sur (sections donc ``de deuxième espèce'') et les composer avec la section donnée. Dans le cas où \deleted{$f$} \inserted{la section rat. $f$ de $U_K$ sur $S$} est \uline{définie} en $s$, il résulte aussitôt de ce qui précède que la section \inserted{partielle} obtenue de $\mathcal E_{U_K}$ au dessus de $\mathcal E_{k(s)}\hookrightarrow\mathcal E_K$ est aussi de première espèce, et \deleted{elle est} (en tant qu'op. \deleted{ext}. \inserted{effective} sur $\pi$) \deleted{\unsure{est}} n'est autre que celle qui est définie par le pt $f(s)\in U_s(k(s))$. Dans
\margin{97}
